\begin{abstract}
Large-scale unsupervised language models (LMs) acquire extensive factual knowledge and display emergent reasoning capabilities; however, directing their outputs remains challenging given the lack of supervision in their original training. Approaches for introducing controllability often rely on human-provided preference data, guiding model outputs via alignment techniques such as reinforcement learning from human feedback (RLHF). Conventional RLHF workflows are intricate and may suffer from instability, typically involving two stages: constructing a reward model from human preferences, then updating the unsupervised LM to optimize for this learned reward, all the while guarding against excessive divergence from the initial model parameters. In this work, we present a novel reward model parameterization within RLHF that permits an analytic solution for the associated optimal policy, thereby enabling resolution of the RLHF objective using a straightforward classification objective. We refer to the resulting method as \textit{Direct Preference Optimization} (DPO); it delivers robust, efficient, and resource-friendly training without the need for language model sampling during the tuning process or laborious hyperparameter adjustments. Through empirical evaluation, we demonstrate that DPO achieves alignment of LMs with human preferences on par with, or superior to, prevalent alternatives. In particular, DPO-based fine-tuning outperforms PPO-style RLHF in steering generative sentiment, and matches or surpasses standard approaches in summarization and dialogue tasks—while offering significant simplifications in implementation and training complexity.
\end{abstract}

\section{Introduction}
Large-scale unsupervised language models (LMs), when trained on massive datasets, exhibit a range of unexpected capabilities~\citep{chowdhery2022palm, brown2020language, touvron2023llama,bubeck2023sparks}. Nevertheless, these models are built on corpora composed of human-generated text, reflecting diverse intentions, skill levels, and objectives. Not all such attributes are beneficial to reproduce; for instance, while it is valuable for an AI-based programming assistant to \textit{recognize} common errors in code for the purpose of correction, during code generation we prefer the model to emulate the (possibly infrequent) examples of superior coding performance embedded in its training set. Analogously, we may seek for our language model to be \textit{informed} about widespread misunderstandings (such as those held by half the population), yet we certainly do not desire the model to assert such inaccuracies as fact in half of the relevant user queries. Thus, discerning and aligning the model's \emph{intended outputs and conduct} with specific desiderata, drawn from its broad \textit{repertoire of information and capabilities}, is a central challenge in the development of safe, reliable, and controllable AI systems~\citep{ouyang2022training}. Traditionally, alignment has been pursued by tuning LMs to reflect human values through reinforcement learning (RL), but as we demonstrate, the standard RL objective underlying these approaches is equivalent to a straightforward binary cross-entropy loss, thereby streamlining the process of preference optimization.

Broadly, prevailing techniques shape language models' behaviors by employing meticulously constructed datasets of human preference annotations, targeting traits regarded as safe and useful. This preference tuning occurs following a preliminary phase of extensive unsupervised pre-training over vast text corpora. The most direct method for integrating preferences involves supervised fine-tuning with high-quality human demonstration data; nonetheless, the leading strategy in current practice relies on reinforcement learning from human (or artificial) feedback (RLHF/RLAIF; \citep{christiano2017deep,bai2022constitutional}). In RLHF, a reward model is first trained on labeled human preference data, after which RL is applied to adjust the model's policy, promoting outputs with high reward scores while constraining deviation from the initial model. Despite the strong performance of RLHF in tasks like conversation and programming, the full RLHF procedure is substantially more intricate than supervised fine-tuning, demanding the training of several models and requiring on-policy sampling throughout optimization, which entails notable computational resources.

This work presents a method for directly aligning a language model with human preferences, eschewing the need for explicit reward modeling or reinforcement learning. We introduce \textit{{Direct Preference Optimization} (DPO)}, a technique that implicitly targets the same optimization criterion as established RLHF approaches—specifically, maximizing reward under a KL-divergence constraint—yet remains highly accessible in terms of implementation and training procedure. The DPO method intuitively augments the log probability assigned to preferred outputs relative to non-preferred ones, integrating a dynamic, sample-specific weighting scheme. This approach mitigates the degradation in model performance associated with naïve probability ratio-based updates. DPO, like traditional techniques, builds on a theoretical foundation of preference modeling (e.g., the Bradley-Terry model; \cite{bradley1952rankanalysis}) for evaluating the correspondence between reward signals and observed preference judgments. However, instead of leveraging the preference model solely for shaping a reward loss that subsequently guides policy optimization, DPO leverages a reparametrization such that the loss is a direct function of the policy itself. Provided a corpus reflecting human choices among model completions, DPO permits direct policy optimization via a straightforward binary cross entropy loss, yielding a policy that is optimal for a latent reward structure inferred from the data.

Our principal innovation is the Direct Preference Optimization (DPO) framework, which furnishes a reinforcement learning-free alternative for preference-based language model fine-tuning. Empirical results indicate that DPO matches or surpasses established preference learning algorithms, including RLHF approaches based on PPO, across tasks such as sentiment control, summarization, and conversational modeling, applied to language models up to 6B parameters in scale.

Large self-supervised language models have demonstrated the capability to address certain tasks in a zero-shot manner \citep{radford2019language} or with the aid of few-shot prompts \citep{gpt3,megatron,chowdhery2022palm}. Nonetheless, their effectiveness on downstream applications and alignment with user intentions are notably enhanced by fine-tuning with datasets composed of instructions and responses authored by humans \citep{mishra-etal-2022-cross,sanh2022multitask,chung2022scaling,thoppilan2022lamda}. This process, referred to as `instruction-tuning,' enables LLMs to handle novel instructions not seen during training, while broadly boosting their overall utility \citep{chung2022scaling}. Although instruction tuning has led to significant improvements, it is often more practical to obtain \textit{relative} assessments of response quality from humans compared to collecting expert demonstrations. Consequently, later work has involved refining LLMs using datasets reflecting human preferences, which has led to gains in tasks such as translation \citep{kreutzer-etal-2018-reliability}, summarization \citep{stiennon2022learning,ziegler2020finetuning}, narrative generation \citep{ziegler2020finetuning}, and instruction adherence \citep{ouyang2022training,ramamurthy2023is}. Typically, these approaches begin by training a neural reward model aligned with the collected preference data under frameworks such as the Bradley-Terry model \citep{bradley1952rankanalysis}. Following this, the language model is further optimized to maximize rewards assigned by this function, generally employing reinforcement learning methods like REINFORCE \citep{williams1992reinforce}, proximal policy optimization (PPO; \cite{schulman2017proximal}), or related techniques \citep{ramamurthy2023is}. Closely allied research has leveraged instruction-tuned LLMs further aligned with human feedback to produce additional synthetic preference datasets, especially targeting criteria such as safety and harmlessness \citep{bai2022constitutional}. These efforts often rely on limited human supervision, such as a text rubric used to guide the LLM’s annotations. This trend reflects the merging of two research directions: one focusing on using reinforcement learning for diverse objectives in language model training~\citep{Ranzato2015SequenceLT,paulus2018a,wu2018learning}, and another on general approaches to preference learning from humans \citep{christiano2017deep,kupcsik2018learning}. Yet, despite the theoretical appeal of leveraging relative human preferences, employing reinforcement learning for fine-tuning large language models introduces substantial practical difficulties; the present work introduces a theoretically-grounded method for optimizing relative preferences without resorting to reinforcement learning.

Beyond the realm of language applications, the problem of learning policies from preference signals has received considerable attention in both bandit and reinforcement learning paradigms, resulting in a variety of methodological contributions. In particular, the study of contextual bandits informed by preferences or ordinal rankings instead of scalar rewards has given rise to the contextual dueling bandit (CDB) framework \citep{yue2012karmed, dudik2015contextual}. In the CDB formulation, since absolute reward information is absent, the concept of an optimal policy is replaced by that of a \textit{von Neumann winner}—a policy that achieves an expected win probability of at least 50\% when compared against any alternative policy \citep{dudik2015contextual}. A salient distinction is that preference information in CDBs is obtained in an online fashion, in contrast to human preference learning scenarios, where training typically proceeds on a fixed, offline dataset of action pairs annotated with preference labels \citep{yan2022human}. Analogously, the \textit{preference-based RL} (PbRL) paradigm learns using binary preferences that originate from an unknown underlying 'scoring' function, eschewing direct reward feedback \citep{BusaFekete2014, ruiz2023dueling}. Various PbRL techniques have been developed, some of which support leveraging off-policy preference datasets. These methods generally require first learning an explicit estimate of the latent scoring or reward function, followed by an optimization step \citep{jain2013learning, BusaFekete2014, christiano2017deep, sadigh2017active, kupcsik2018learning}. In contrast, the approach we propose integrates policy optimization and preference satisfaction into a single stage, circumventing the need for a separate reward modelling step.

\section{Preliminaries}\label{section:prelims}

We briefly summarize the RLHF workflow introduced by \citeauthor{ziegler2020finetuning} and subsequently adopted by others \citep{stiennon2022learning, bai2022training, ouyang2022training}. This pipeline generally involves three main stages: 1) supervised fine-tuning (SFT); 2) gathering preference data and learning a reward model; and 3) reinforcement learning-based policy improvement.

\textbf{SFT}: The RLHF process typically commences with supervised fine-tuning of a pre-trained language model, using a curated corpus suited to the downstream task of interest (e.g., dialogue or summarization), thereby yielding the model $\pi^\text{SFT}$.

\textbf{Reward Modelling Phase}: In the subsequent phase, the fine-tuned SFT model receives prompt inputs $x$ and generates pairs of outputs $(y_1, y_2)\sim \pi^\text{SFT}(y \mid x)$. Human annotators are then presented with these answer pairs and express a preference, denoted $y_w\succ y_l \mid x$, with $y_w$ and $y_l$ indicating the more and less preferred answers, respectively. The underlying assumption is that preferences are governed by an unobserved reward function $r^*(y, x)$, which is not directly accessible. Several modeling approaches exist for capturing these preference signals; the Bradley-Terry (BT) model \cite{bradley1952rankanalysis} is commonly adopted (although the more general Plackett-Luce ranking model \citep{plackett1975analysis, luce2012individual} can also be utilized when multiple responses are ranked). In the BT framework, the human preference distribution $p^*$ is defined as:

Given access to a fixed dataset of pairwise comparisons $\mathcal{D}=\bigl\{x^{(i)}, y_w^{(i)}, y_l^{(i)}\bigr\}_{i=1}^N$ drawn according to $p^*$, one can introduce a reward model $r_{\phi}(x, y)$ parameterized by $\phi$ and learn its parameters through maximum likelihood estimation. If we frame the task as binary classification, the resulting objective corresponds to the negative log-likelihood:

\begin{equation}\label{eq:reward_model}
    \mathcal{L}_R(r_{\phi}, \mathcal{D}) = -\mathbb{E}_{(x, y_w, y_l)\sim \mathcal{D}}\left[\log \sigma\left(r_{\phi}(x, y_w) - r_{\phi}(x, y_l)\right)\right]
\end{equation}

where $\sigma$ is the logistic sigmoid. For language models, it is standard practice to initialize $r_{\phi}(x, y)$ from the SFT model $\pi^\text{SFT}(y \mid x)$, appending a linear projection atop the final transformer layer, yielding a scalar reward output as in \cite{ziegler2020finetuning}. Previous studies additionally enforce reward normalization, often setting  $\mathbb{E}_{x,y\sim \mathcal{D}}\left[r_\phi(x, y)\right]=0$ per context $x$, in order to stabilize the reward function and minimize its variance.

\textbf{RL Fine-Tuning Phase}: In the subsequent reinforcement learning step, the reward function is deployed to provide evaluative signals for policy improvement. Building upon approaches from earlier work~\citep{jaques2017sequence, jaques2020human}, the optimization objective takes the following form:

\begin{equation}\label{eq:RL}
\max_{\pi_{\theta}} \mathbb{E}_{x\sim \mathcal{D}, y\sim \pi_{\theta}(y \mid x)}\left[r_{\phi}(x, y)\right] - \beta \mathbb{D}_{\textrm{KL}}\left[\pi_{\theta}(y\mid x)\,||\,\pi_\text{ref}(y\mid x)\right],
\end{equation}

where $\beta$ serves as a regularization coefficient that restricts $\pi_\theta$ from straying too far from the reference distribution $\pi_\text{ref}$, which is usually chosen to be the pre-trained SFT policy $\pi^\text{SFT}$. Conventionally, the policy $\pi_\theta$ is initialized to this same SFT model. Imposing this KL divergence penalty is crucial for preventing excessive divergence from regions of the input space where the reward model remains reliable, as well as preserving sample diversity and reducing the risk of collapse to narrow high-reward modes. Since language generation is inherently discrete, this objective is non-differentiable and thus reinforcement learning methods are used. In practice, a composite reward function ${r(x, y) = r_{\phi}(x, y) -\beta (\log \pi_{\theta}(y\mid x) - \log \pi_\text{ref}(y\mid x))}$ is constructed and optimized using PPO \cite{schulman2017proximal}, following prior work \citep{ziegler2020finetuning, stiennon2022learning, bai2022training, ouyang2022training}.

\section{Direct Preference Optimization}\label{sec:DPO}

Driven by the considerable complexity inherent in applying reinforcement learning algorithms to the large-scale fine-tuning of language models, we aim to establish a more straightforward technique that directly leverages preference data for policy optimization. In contrast to traditional RLHF methods, which first fit a reward model before optimizing it via reinforcement learning, our framework utilizes a specific reward parameterization that admits a closed-form expression for its optimal policy—bypassing the need for an iterative RL loop. As we elaborate below, our central contribution is an analytic transformation from reward models to optimal policy distributions. This “change of variables” reframes loss functions on rewards into loss functions on policies, thereby enabling direct policy optimization aligned with human preferences models such as the Bradley-Terry model—without explicitly fitting an intermediate reward predictor. Thus, in our formulation, the policy model is trained to implicitly encode both the generation and the reward estimation.

\textbf{Derivation of the DPO objective.} Beginning with the same reinforcement learning objective as used in previous literature (see Eq.~\ref{eq:RL}), and considering a generic reward function $r$, it is well established~\citep{peters2007reinforcement, peng2019advantage, korbak2022reinforcement, go2023aligning} that the KL-regularized reward maximization objective given in Eq.~\ref{eq:RL} yields an optimal policy of the following form:

\begin{equation}\label{eq:op_policy}
    \pi_r(y\mid x) = \frac{1}{Z(x)}\pi_\text{ref}(y\mid x)\exp\left(\frac{1}{\beta}r(x, y)\right),
\end{equation}

Here, the partition function $Z(x) =\sum_{y}\pi_\text{ref}(y\mid x)\exp\left(\frac{1}{\beta}r(x, y)\right)$ ensures normalization. The complete derivation is available in Appendix \ref{app:derivation1}. Even substituting in the maximum likelihood estimate $r_\phi$ for the true reward $r^*$, evaluating $Z(x)$ remains computationally demanding~\citep{korbak2022reinforcement, go2023aligning}, rendering this formulation less tractable for practical use. To circumvent this, one can rearrange Eq.~\ref{eq:op_policy} so that the reward is expressed using its optimal policy $\pi_r$, reference policy $\pi_\text{ref}$, and the unobserved $Z(x)$. By applying the logarithm to both sides of Eq.~\ref{eq:op_policy}, some straightforward algebra yields:

\begin{equation}\label{eq:main_eq}
    r(x,y) =\beta \log \frac{\pi_r(y\mid x)}{\pi_\text{ref}(y\mid x)} + \beta \log Z(x).
\end{equation}

We apply this change of variables to both the actual reward $r^*$ and its corresponding optimal policy $\pi^*$. Importantly, within the Bradley-Terry framework, the human preference between two completions only depends on the difference of their respective rewards, i.e., ${p^*(y_1 \succ y_2 \mid x) = \sigma(r^*(x, y_1) - r^*(x, y_2))}$. Replacing $r^*(x, y)$ in the Bradley-Terry model Eq.~\ref{eq:bradley-terry} using the parametrization from Eq.~\ref{eq:main_eq} leads to the cancellation of the partition function, allowing the preference probability to be written purely in terms of the optimal and reference policies. Thus, under the Bradley-Terry assumption, the optimal RLHF policy $\pi^*$ must satisfy:

\begin{equation}\label{eq:objective}
    p^*(y_1\succ y_2 \mid x)=\frac{1}{1 + \exp\left(\beta \log \frac{\pi^*(y_2\mid x)}{\pi_\text{ref}(y_2\mid x)} - \beta \log \frac{\pi^*(y_1\mid x)}{\pi_\text{ref}(y_1\mid x)}\right)}
\end{equation}

Detailed steps are in Appendix~\ref{app:derivation2}. While this derivation is shown for the Bradley-Terry model, a parallel development can be made for the broader class of Plackett-Luce models~\citep{plackett1975analysis, luce2012individual}, with details in Appendix~\ref{app:plackett_luce_models}.

Given this expression for preference probability in terms of the optimal policy (as opposed to the reward function), it is now possible to specify a maximum likelihood objective for a parameterized policy $\pi_\theta$. Analogous to the reward modeling objective (e.g., Eq.~\ref{eq:reward_model}), we have:

\begin{equation}\label{eq:optimum_model}
    \mathcal{L}_\text{DPO}(\pi_{\theta}; \pi_\text{ref}) = -\mathbb{E}_{(x, y_w, y_l)\sim \mathcal{D}}\left[\log \sigma \left(\beta \log \frac{\pi_{\theta}(y_w\mid x)}{\pi_\text{ref}(y_w\mid x)} - \beta \log \frac{\pi_{\theta}(y_l\mid x)}{\pi_\text{ref

Through this approach, we model an implicit reward by adopting an alternative parameterization, in which the optimal policy directly corresponds to $\pi_\theta$. Furthermore, as our methodology parallels the process of fitting a reparameterized Bradley-Terry model, it inherits beneficial theoretical guarantees—most notably, consistency under appropriate conditions on the preference data distribution \cite{bong2022generalized}. Section~\ref{sec:theory} offers a deeper exploration of these theoretical aspects of DPO, situating them within the context of existing literature.

\textbf{Mechanistic Interpretation of the DPO Update.} To gain a mechanistic perspective on the DPO algorithm, it is instructive to examine the gradient of its loss function $\mathcal{L}_\text{DPO}$. The gradient with respect to the parameter vector $\theta$ can be expressed as:

\begin{multline*}\label{eq:gradient}
    \nabla_\theta \mathcal{L}_\text{DPO}(\pi_\theta;\pi_\text{ref}) = \\ -\beta\mathbb{E}_{(x, y_w, y_l) \sim \mathcal{D}} \bigg[\underbrace{\sigma(\hat{r}_\theta(x, y_l) - \hat{r}_\theta (x, y_w))}_\text{assigns higher weights to misranked rewards}\bigg[\underbrace{\nabla_\theta\log \pi(y_w \mid x)}_\text{promotes higher likelihood for $y_w$} - \underbrace{\nabla_\theta\log\pi(y_l \mid x)}_\text{reduces likelihood for $y_l$}\bigg]\bigg],
\end{multline*}

In this context, $\hat{r}_\theta(x, y) = \beta \log \frac{\pi_\theta(y \mid x)}{\pi_\text{ref}(y \mid x)}$ represents the implicit reward inferred by the language model $\pi_\theta$ relative to the reference policy $\pi_\text{ref}$ (see also Section~\ref{sec:theory}). At a high level, the gradient of $\mathcal{L}_\text{DPO}$ is designed to reinforce the probability assigned to preferred responses $y_w$, while simultaneously diminishing the probability of less preferred responses $y_l$. Crucially, each training sample is weighted by the extent to which the implicit reward model $\hat{r}_\theta$ erroneously ranks the dispreferred completion above the preferred one, modulated by the parameter $\beta$, thereby incorporating both the degree of misranking and the magnitude of the KL divergence constraint. Empirically, this form of weighting proves essential: omitting it and applying a uniform weighting—i.e., removing the weighting coefficient—can result in pathological behavior and model degeneration (see Appendix Table~\ref{tab:unlikelihood_generations}).

\textbf{DPO overview.} The standard workflow for DPO consists of the following steps: 1) For each prompt $x$, generate two candidate completions $y_1, y_2 \sim \pi_\text{ref}(\cdot \mid x)$, assign preference labels via human annotation, and thereby form an offline preference dataset $\mathcal{D} = \{x^{(i)}, y_w^{(i)}, y_l^{(i)}\}_{i=1}^N$; 2) update the language model $\pi_\theta$ by minimizing $\mathcal{L}_\text{DPO}$, conditioned on the fixed $\pi_\text{ref}$, dataset $\mathcal{D}$, and a chosen $\beta$ parameter.

In applied settings, researchers often favor leveraging existing public preference datasets over assembling new data via sampling and human annotation. Typically, such datasets are constructed using samples from $\pi^\text{SFT}$. Thus, whenever this model is accessible, we set $\pi_\text{ref} = \pi^\text{SFT}$. In cases where $\pi^\text{SFT}$ is unavailable, $\pi_\text{ref}$ is instead initialized by maximizing the likelihood of preferred completions ${(x, y_w)}$ in the dataset, specifically, by solving ${\pi_\text{ref} = \argmax_{\pi}\mathbb{E}_{x, y_w \sim \mathcal{D}}\left[\log \pi(y_w \mid x)\right]}$. This strategy serves to alleviate the mismatch between the unavailable true reference distribution and the practical choice of $\pi_\text{ref}$ in DPO. Comprehensive information on implementation choices and hyperparameter selection is provided in Appendix~\ref{app:implementation}.

\section{Theoretical Analysis of DPO} In this section, we further interpret the DPO approach, furnish its theoretical foundation, and discuss its benefits in comparison to actor-critic methods common in RLHF (notably PPO~\cite{schulman2017proximal}).

\label{sec:theory}

\subsection{Your Language Model Is Secretly a Reward Model} DPO uniquely avoids explicit reward modeling and reinforcement learning steps for policy optimization, unifying these steps through a single maximum likelihood objective. In particular, the DPO objective (Eq. \ref{eq:main_eq}) can be reformulated as a Bradley-Terry model, where the reward is parameterized as $r^*(x, y) = \beta \log\frac{\pi^*_\theta(y \mid x)}{\pi_\text{ref}(y \mid x)}$, and we optimize the parametric model $\pi_{\theta}$—analogous, after a variable change, to optimizing a reward model as in Eq. \ref{eq:reward_model}. This section develops the theoretical perspective underlying this reparameterization, demonstrating that it imposes no restrictions on the expressiveness of the learned reward functions and permits the exact identification of the optimal policy. We start by establishing an equivalence relation on reward functions.

\begin{definition}
Two reward functions $r(x, y)$ and $r'(x, y)$ are considered equivalent if and only if there exists some function $f$ such that ${r(x, y)-r'(x, y) = f(x)}$.
\end{definition}

This equivalence relation is straightforward to verify and results in the partitioning of reward functions into distinct equivalence classes. We now present two key lemmas:

\begin{lemma}\label{lemma:same_prefrence}
Within the Plackett-Luce preference framework—including the special case of Bradley-Terry—any pair of reward functions belonging to the same equivalence class will generate identical preference distributions.
\end{lemma}

\begin{lemma}\label{lemma:same_policy}
    Any two reward functions drawn from the same equivalence class produce identical optimal policies when considering the constrained RL problem.
\end{lemma}

The demonstrations of these statements are uncomplicated and are provided in Appendix \ref{app:lemma1}. The first lemma captures a recognized under-specification phenomenon characteristic of the Plackett-Luce model family \cite{plackett1975analysis}. This ambiguity necessitates the introduction of identifiability constraints if one wishes to obtain well-defined MLE solutions, as referenced in Eq. \ref{eq:reward_model} \cite{bong2022generalized}. The second lemma asserts that all reward functions in an equivalence class lead to the same optimal policy, implying that for policy recovery, it suffices to recover any representative reward function from the correct equivalence class. In Appendix~\ref{app:thm1}, we provide the proof of the following theorem:

\begin{theorem}\label{thm:main}
    Provided mild regularity conditions hold, every reward class compatible with the Plackett-Luce family (specifically, including Bradley-Terry) can be parameterized as ${r(x, y) = \beta \log \frac{\pi(y\mid x)}{\pi_\text{ref}(y\mid x)}}$, for some model $\pi(y\mid x)$ and a fixed reference distribution $\pi_\text{ref}(y \mid x)$.
\end{theorem}

\begin{sproof}
    Take any given reward function $r(x, y)$ and the associated optimal model $\pi_r(y \mid x)$, as defined in Eq. \ref{eq:op_policy}. Our aim is to show that there exists a member in the equivalence class of $r$ that can be represented by the above reparameterization. We construct the mapping $f$ by
\begin{equation}
    f(r; \pi_\text{ref}, \beta)(x, y) = r(x, y) - \beta\log\sum_{y}\pi_\text{ref}(y\mid x)\exp\left(\frac{1}{\beta}r(x, y)\right)
\end{equation}
This mapping $f$ normalizes the reward by subtracting the logarithm of the corresponding partition function (which depends only on $x$), ensuring $f(r; \pi_\text{ref}, \beta)(x, y)$ remains within the equivalence class of $r(x, y)$. By substituting $r$ with the right side of Eq.~\ref{eq:main_eq} (which applies for all reward functions), we see that $f(r; \pi_\text{ref}, \beta)(x, y) = \beta \log \frac{\pi_r(y\mid x)}{\pi_\text{ref}(y\mid x)}$. Thus, $f$ yields an equivalence-class representative of $r$ in the stated canonical form, ensuring no loss of generality with this reparameterization.
\end{sproof}

An alternative interpretation of Theorem~\ref{thm:main} is that it identifies the precise reward function, within each equivalence class, that is chosen by the DPO reparameterization. Specifically, it selects the reward function $r(x, y)$ that fulfills:

\begin{equation}\label{eq:lag_p}
     \sum_{y}\underbrace{\pi_\text{ref}(y\mid x)\exp\left(\frac{1}{\beta}r(x, y)\right)}_{=\pi(y\mid x)\text{, using Thm.~\ref{thm:main} reparam.}} = 1,
\end{equation}

meaning that $\pi(y\mid x)$ constitutes a proper probability distribution (i.e., all probabilities are non-negative and sum to one). Moreover, when considered in light of Eq.~\ref{eq:op_policy}, Eq.~\ref{eq:lag_p} corresponds to the normalization constant—often referred to as the partition function—of the optimal policy determined by $r(x, y)$. The central insight underpinning the DPO method is the ability to introduce specific constraints into the underdetermined Plackett-Luce (including, as a notable instance, the Bradley-Terry) class of preference models. This approach retains the scope of representable reward functions while simultaneously ensuring that the optimal policy in Eq.~\ref{eq:op_policy} is analytically computable for any prompt $x$.

\subsection{Instability of Actor-Critic Algorithms}
Our analytical framework also serves as a tool to better understand sources of instability in conventional actor-critic methods commonly used for RLHF, such as PPO. Focusing on the RL fine-tuning phase, as detailed in Section \ref{section:prelims}, we can draw parallels with the control-as-inference paradigm~\cite{levine2018reinforcement}, applied here to the constrained RL setup introduced in \ref{eq:RL}. Suppose we have a model $\pi_{\theta}(y\mid x)$ parameterized by $\theta$; the objective is to minimize the KL divergence $\mathbb{D}_{\text{KL}}[\pi_{\theta}(y|x) \mid \mid \pi^*(y\mid x)]$, where the reference optimal policy $\pi^*$, derived from Eq.~\ref{eq:optimum_model}, depends on the reward function $r_{\phi}(y, x)$. By expanding this formulation, we obtain the following optimization criterion:

\begin{equation}\label{eq:AC}
    \max_{\pi_{\theta}}\mathbb{E}_{\pi_{\theta}(y\mid x)}\left[\underbrace{r_{\phi}(x, y) -\beta\log\sum_{y}\pi_\text{ref}(y\mid x)\exp\left(\frac{1}{\beta}r_{\phi}(x, y)\right)}_{f(r_{\phi}, \pi_\text{ref}, \beta)} - \underbrace{\beta\log\frac{\pi_{\theta}(y\mid x)}{\pi_\text{ref}(y\mid x)}}_{\text{KL}}\right]
\end{equation}

The objective here is identical to that optimized in earlier research 
\citep{ziegler2020finetuning, stiennon2022learning, bai2022training, ouyang2022training}, where the DPO-equivalent reward is employed for the $r_{\phi}$ reward category. In this framework, the normalization component within $f(r_{\phi}, \pi_\text{ref}, \beta)$ can be understood as the soft value function associated with the reference policy $\pi_\text{ref}$. Although this normalization factor does not alter the optimal policy, omitting it may result in a policy gradient with high variance, thereby reducing the stability of the training process. One can address the normalization through a learned value function, though this presents optimization challenges as well. Alternatively, previous work has normalized rewards based on a baseline derived from human completions, which effectively corresponds to a single-sample Monte Carlo estimation of the normalization constant. The DPO reparameterization, in comparison, yields a reward function that eliminates the necessity for such baselines.

\section{Experiments}
This section presents a series of empirical analyses to assess the capacity of DPO to train policies directly from preference data. We first investigate, under controlled conditions for text generation, the efficiency with which DPO balances reward maximization and KL-divergence minimization relative to the reference policy, benchmarking against commonly used preference learning algorithms like PPO. We then examine DPO's effectiveness on larger-scale models and more complex RLHF applications, specifically focusing on summarization and dialogue tasks. Our findings indicate that, with minimal hyperparameter adjustment, DPO generally matches or outperforms robust baselines, including RLHF with PPO, and outperforms strategies such as selecting the best among $N$ trajectories according to a learned reward function. The details of our experimental procedure are provided below, with further specifics available in Appendix~\ref{app:exp_details}.

\textbf{Tasks.} We investigate three distinct open-ended text generation tasks in our experiments. Across all tasks, the algorithms are trained to optimize a policy using a dataset of preference pairs $\mathcal{D}=\bigl\{x^{(i)}, y_w^{(i)}, y_l^{(i)}\bigr\}_{i=1}^N$. In the case of \textbf{controlled sentiment generation}, $x$ serves as the beginning segment of a movie review extracted from the IMDb dataset \cite{maas-EtAl:2011:ACL-HLT2011}, with the objective that the generated continuation $y$ expresses positive sentiment. To facilitate controlled assessment, we automatically construct preference pairs for generations by leveraging a pre-trained sentiment classifier, ensuring that $p(\text{positive}\mid x,y_w)>p(\text{positive}\mid x,y_l)$. For SFT, GPT-2-large is further trained to convergence on the IMDB dataset’s training partition (additional information in App~\ref{app:sentiment_details}). The \textbf{summarization} task uses $x$ as a Reddit forum post; here, the policy aims to produce a summary $y$ capturing the core ideas. As in previous studies, we utilize the Reddit TL;DR dataset \citep{volske-etal-2017-tl} as well as human preference annotations from \citeauthor{stiennon2022learning}. The SFT model is obtained by fine-tuning on summaries crafted by humans\footnote{\url{https://huggingface.co/CarperAI/openai_summarize_tldr_sft}}, employing the TRLX framework \citep{leandro_von_werra_2023_7790115} for RLHF. The collection of human preferences, used for evaluation, was assembled by \citeauthor{stiennon2022learning} from outputs of a related SFT model. In the third setting, \textbf{single-turn dialogue}, $x$ denotes a human input that may span topics from astrophysics queries to requests for advice. The policy is expected to generate a helpful and engaging reply $y$; we utilize the Anthropic Helpful and Harmless dialogue dataset \citep{bai2022training}, featuring 170k conversational exchanges between humans and an automated system. Each dialogue ends with two alternative model-generated replies and an accompanying annotation indicating which is preferred by a human judge. As a pre-existing SFT model is unavailable for this dataset, we construct one by fine-tuning a generic language model solely on the preferred outputs.

\textbf{Evaluation.} We assess performance through two distinct evaluation strategies. In the controlled sentiment generation scenario, where the ground-truth reward function (specifically, a sentiment classifier) is accessible, we systematically measure each algorithm by plotting the trade-off frontier between the realized reward and KL-divergence relative to the reference policy, thereby allowing direct examination of reward maximization under constraints. Conversely, practical applications typically lack an explicit reward function; thus, we shift to a \textit{win rate} evaluation paradigm. Here, we benchmark models against a baseline policy, leveraging GPT-4 as an automated proxy for human judgment to assess the quality of summaries and the helpfulness of single-turn dialogue responses. For summarization tasks, the test set’s reference summaries serve as baselines, while for dialogue, baseline comparisons use the human-preferred response provided in the dataset. While literature supports the reliability of language models as evaluators over traditional metrics \citep{Chen2023ExploringTU}, we complement our approach with a human subject study—detailed in Sec.~\ref{sec:human-judgments}—which confirms the validity of our methodology: agreement between human raters and GPT-4 is high, and often matches or surpasses human-to-human concordance.

\textbf{Methods.} Beyond DPO, our analysis encompasses multiple existing strategies for aligning language models with human preferences. For the summarization domain, we implement zero-shot prompting with \textbf{GPT-J} \citep{gpt-j}; in dialogue, we use 2-shot prompting with \textbf{Pythia-2.8B} \citep{biderman2023pythia}. We further evaluate the \textbf{SFT} approach, alongside \textbf{Preferred-FT}, wherein supervised fine-tuning is performed on completions $y_w$ selected either from SFT outputs (for sentiment control and summarization) or generated by a generic language model (for dialogue tasks). Another supervised-style technique is the \textbf{Unlikelihood} method~\citep{welleck2019neural}, which encourages the model to increase probability mass on $y_w$ while suppressing it on $y_l$, modulated by an optional `unlikelihood' weight $\alpha\in[0,1]$. Reinforcement learning approaches under consideration include \textbf{PPO} \citep{schulman2017proximal}—using a reward model trained from preferences—and \textbf{PPO-GT}, an oracle variant informed by the available ground-truth reward in sentiment-controlled tasks. For sentiment analysis, we implement two PPO-GT versions: one as provided in \cite{leandro_von_werra_2023_7790115}, and another modified variant with normalized rewards and refined hyperparameters, improvements that are similarly applied to PPO with learned rewards. Finally, we introduce a strong but computationally demanding \textbf{Best of $N$} baseline, which generates $N$ samples from the SFT model (or Preferred-FT for dialogue) and returns the response rated highest by the preference-trained reward model. This method isolates reward model quality from PPO optimization but entails sampling $N$ completions for each test instance, rendering it inefficient for even moderate $N$.

\subsection{How well can DPO optimize the RLHF objective?}

The KL-constrained reward maximization objective, fundamental to most RLHF algorithms, seeks to optimize for reward while imposing a penalty on divergence from a reference policy. Consequently, when evaluating different methods, it is essential to consider both the achieved reward and the magnitude of KL divergence; a marginal increase in reward accompanied by a substantial rise in KL is often undesirable. Figure~\ref{fig:frontier-tldr-main} illustrates the relationship between reward and KL for several algorithms within the sentiment domain. For each method, we perform several training sessions, systematically varying a hyperparameter associated with policy conservativeness: the target KL ($\in\{3,6,9,12\}$) for PPO, $\beta$ ($\in \{0.05,0.1,1,5\}$), and $\alpha$ ($\in\{0.05,0.1,0.5,1\}$) for the unlikelihood objective, along with different random seeds for preferred-FT, resulting in a total of 22 experiment runs. At every 100 training steps up to convergence, we evaluate the current policy using a set of test prompts, recording both the mean reward under the true reward function and the mean sequence-level KL\footnote{Defined as the aggregate of the KL-divergence at each timestep.} with respect to the reference policy $\text{KL}\left(\pi\mid \mid \pi_\text{ref}\right)$. Our results indicate that DPO delivers the most favorable frontier, simultaneously achieving high reward and maintaining a low KL. This finding stands out for several reasons. Notably, despite both DPO and PPO targeting the same objective, DPO consistently exhibits greater efficiency; its reward-to-KL profile entirely surpasses that of PPO. Furthermore, DPO attains a superior frontier compared to PPO even when PPO is provided with access to the ground truth reward signal (PPO-GT).

\subsection{Can DPO scale to real preference datasets?}
\label{sec:dpo-real-datasets}
We proceed to assess how effectively DPO fine-tunes models on real-world tasks such as summarization and single-turn dialogue. In the domain of summarization, metrics like ROUGE are often found to correlate weakly with human judgment~\citep{stiennon2022learning}, and earlier studies have shown that leveraging PPO to fine-tune language models according to human preferences yields more satisfactory summaries. Our evaluation methodology involves generating responses from each method using the test partition of the TL;DR summarization dataset and measuring the average win rate relative to the provided reference summaries within the test set. For all approaches, responses are sampled across a range of temperatures between 0.0 and 1.0, with results presented in Figure~\ref{fig:frontier-tldr-main} (right). DPO, PPO, and Preferred-FT all build upon the identical GPT-J SFT model as a starting point\footnote{\url{https://huggingface.co/CarperAI/openai_summarize_tldr_sft}}. Empirically, DPO achieves an approximate win rate of 61\% when sampled at temperature 0.0, surpassing PPO's best win rate of roughly 57\% at the same temperature. Furthermore, DPO's peak win rate exceeds that of the strongest $N$ baseline. It is important to highlight that we performed minimal tuning of DPO’s $\beta$ parameter, implying that our reported numbers may underrepresent DPO's full capability. Additionally, we observe that DPO exhibits considerably greater stability with respect to sampling temperature changes compared to PPO, whose win rate drops to match that of the base GPT-J model as temperature increases. Meanwhile, Preferred-FT displays only marginal improvement relative to the SFT baseline. A direct comparison between DPO and PPO using human raters is presented in Section~\ref{sec:human-judgments}, where samples from DPO at temperature 0.25 are favored 58\% of the time over PPO outputs at the same temperature.

For single-turn dialogue, we assess various approaches using the subset of the Anthropic HH dataset \citep{bai2022training} test split that features a single round of interaction between a human and assistant. In the GPT-4 assessment, we utilize the preferred completions from the test set as references, measuring each method's win rate against these. Due to the absence of a standardized SFT baseline for this evaluation, we begin with a pre-trained Pythia-2.8B, employ Preferred-FT to develop a reference model by training it on selected completions so that the outputs remain in-distribution, and then apply DPO for further training. Our comparisons also include the top result from 128 Preferred-FT completions (noting that the Best of $N$ approach plateaus with $N=128$ in this context; refer to Appendix Figure~\ref{fig:best-of-n}) and a 2-shot prompting setup for the Pythia-2.8B base model. Across temperature settings, we observe DPO matching or exceeding the performance of the best-tuned alternatives. Furthermore, we evaluate an RLHF model fine-tuned with PPO on the Anthropic HH dataset\footnote{\url{https://huggingface.co/reciprocate/ppo_hh_pythia-6B}} sourced from a reputable repository\footnote{\url{https://github.com/CarperAI/trlx/tree/main/examples/hh}}, but we are unable to identify a prompt or sampling temperature that results in better outcomes than the unadapted Pythia-2.8B model. Drawing on our TL;DR results and recognizing that both Best of 128 and PPO optimize equivalent reward functions, we treat the Best of 128 method as an approximate benchmark for PPO performance. Notably, DPO stands out as the sole approach that both efficiently utilizes computation and yields improvements over the preferred completions in the Anthropic HH dataset, matching or surpassing the results of the more resource-intensive Best of 128 baseline. Lastly, Figure~\ref{fig:dialogue-main} illustrates that DPO achieves its optimal performance in relatively few training iterations.

\subsection{Generalization to a new input distribution}

To more thoroughly assess how PPO and DPO handle distribution shifts, we apply the policies derived from our Reddit TL;DR summarization experiments to a distinct dataset: news articles from the CNN/DailyMail test set \citep{nallapati-etal-2016-abstractive}, using the optimal sampling temperatures determined for TL;DR (0 and 0.25). The outcomes are detailed in Table~\ref{tab:ood}. We measured the rate at which GPT-4 preferred the model-generated summaries over the reference summaries, adapting the same GPT-4 (C) prompt from the Reddit TL;DR experiments by replacing ``forum post'' with ``news article.'' On this out-of-distribution test, DPO significantly surpasses the PPO policy. These findings provide preliminary support that DPO policies exhibit generalization capabilities comparable to those of PPO policies, despite DPO not leveraging the supplementary unlabeled Reddit TL;DR prompts that inform PPO training.

\subsection{Validating GPT-4 judgments with human judgments}
\label{sec:human-judgments}
We implement a human evaluation to assess the validity of GPT-4-based preference judgments, using the TL;DR summarization experiment outcomes and two alternative GPT-4 prompts. The \textbf{GPT-4 (S)} (simple) prompt requests a judgment on which summary better encapsulates the essential content of the post. In contrast, the \textbf{GPT-4 (C)} (concise) prompt inquires which summary is both more informative and concise; this is motivated by observations that, under the \textbf{GPT-4 (S)} prompt, GPT-4 tends to favor lengthier and more repetitive summaries compared to human raters. Complete prompt texts can be found in Appendix~\ref{app:prompts}. We conduct three sets of pairwise comparisons—between the top-performing (DPO, temp. 0.25), the lowest-performing (PPO, temp. 1.0), and a mid-range method (SFT, temp. 0.25)—contrasted against PPO sampled greedily at its best temperature. This design captures a spectrum of output qualities. Analysis indicates that for both prompt variants, GPT-4’s agreement with human evaluations is on par with human-human agreement rates, implying GPT-4 is a suitable stand-in for direct human assessment (with multiple human ratings collected only for DPO and PPO-1, due to a limited rater pool). Among the two, the \textbf{GPT-4 (C)} prompt produces win rates more aligned with human preferences, and thus is used for our principal results in Section~\ref{sec:dpo-real-datasets}. Appendix~\ref{app:human-study} contains further information on the study methodology, the interface provided to raters, and the list of human evaluators.

\section{Discussion}
Preference-based learning offers a robust and scalable approach for developing proficient and aligned language models. In this work, we introduced DPO, a straightforward methodology for leveraging human preferences in language model training that circumvents the need for reinforcement learning. Instead of adapting the preference learning task into a traditional RL framework for compatibility with existing RL solutions, DPO formulates a direct relationship between language model policies and associated reward functions. This enables the training process to align with human feedback using a simple cross-entropy loss, eschewing reinforcement learning entirely without compromising on generality. Remarkably, DPO achieves performance on par with, or superior to, popular RLHF methods—including PPO-based strategies—while demanding minimal hyperparameter adjustment. As a result, DPO notably simplifies the process of preference-based language model training.

\textbf{Limitations \& Future Work.} This study prompts a number of compelling avenues for continued research. For instance, how does the generalization ability of policies trained via DPO compare to those trained on explicit reward functions, especially outside of the training distribution? Preliminary findings indicate comparable generalization between DPO and PPO approaches, yet this requires more thorough analysis. Another open question concerns the effectiveness of self-labeling—can DPO-trained policies fully leverage unannotated prompts through this mechanism? Moreover, the phenomenon of reward over-optimization warrants investigation in the context of direct preference optimization; it remains to be seen if the minor dip in performance observed in Figure~\ref{fig:dialogue-main}-right is a reflection of this. While our current experiments utilize models up to 6B parameters, scaling DPO for application to much larger, cutting-edge models presents a promising research trajectory. Additionally, evaluation procedures deserve further scrutiny, as our experiments reveal that GPT-4-derived win rates are sensitive to prompt phrasing; future investigations might seek optimal practices for eliciting reliable judgments from automated evaluators. Lastly, DPO's utility may well extend beyond aligning language models with human preferences, with the potential for application in the broader domain of generative modeling across various data modalities.